\documentclass[12pt,a4paper]{article}

% Paquetes
\usepackage[spanish]{babel}
\usepackage[utf8]{inputenc}
\usepackage[T1]{fontenc}
\usepackage{geometry}
\usepackage{hyperref}
\usepackage{graphicx}
\usepackage{xcolor}
\usepackage{booktabs}
\usepackage{enumitem}
\usepackage{fancyhdr}
\usepackage{titlesec}
\usepackage{mdframed}
\usepackage{lmodern}
\usepackage{microtype}
\usepackage{parskip}

% Geometría
\geometry{
  a4paper,
  top=2.5cm, bottom=2.5cm,
  left=2.8cm, right=2.8cm
}

% Colores
\definecolor{primaryColor}{RGB}{99,102,241}   % índigo
\definecolor{successColor}{RGB}{16,185,129}   % verde
\definecolor{warningColor}{RGB}{245,158,11}   % ámbar
\definecolor{dangerColor}{RGB}{239,68,68}     % rojo
\definecolor{grayLight}{RGB}{243,244,246}
\definecolor{grayText}{RGB}{107,114,128}

% Encabezados de sección
\titleformat{\section}
  {\large\bfseries\color{primaryColor}}
  {\thesection.}{0.5em}{}[\titlerule]

\titleformat{\subsection}
  {\normalsize\bfseries}
  {\thesubsection.}{0.5em}{}

% Cabecera y pie de página
\pagestyle{fancy}
\fancyhf{}
\fancyhead[L]{\small\textcolor{grayText}{Manual del Panel Administrativo — Corazón Migrante}}
\fancyhead[R]{\small\textcolor{grayText}{\thepage}}
\renewcommand{\headrulewidth}{0.4pt}

% Recuadros de consejo
\newenvironment{consejo}{%
  \begin{mdframed}[backgroundcolor=grayLight,linecolor=primaryColor,linewidth=1.5pt,innerleftmargin=10pt]
  \textcolor{primaryColor}{\textbf{Consejo:}} 
}{%
  \end{mdframed}
}

\newenvironment{importante}{%
  \begin{mdframed}[backgroundcolor=grayLight,linecolor=dangerColor,linewidth=1.5pt,innerleftmargin=10pt]
  \textcolor{dangerColor}{\textbf{Importante:}} 
}{%
  \end{mdframed}
}

% ================================================================
\begin{document}

% Portada
\begin{titlepage}
  \centering
  \vspace*{3cm}
  {\huge\bfseries\color{primaryColor} Manual del Panel Administrativo\\[0.5em]}
  {\large Corazón Migrante\\[2em]}
  {\normalsize Para psicólogos y coordinadores de la plataforma\\[0.3em]}
  {\small Sin conocimientos técnicos necesarios\\[4cm]}
  \hrule
  \vspace{0.5em}
  {\small Versión 1.0 \textperiodcentered{} Julio 2026}
  \vspace{0.5em}
  \hrule
  \vfill
  {\footnotesize\textcolor{grayText}{Este documento es de uso interno. No distribuir sin autorización.}}
\end{titlepage}

% Índice
\tableofcontents
\newpage

% ================================================================
\section{Introducción}

Este manual explica cómo usar el panel administrativo de Corazón Migrante.
Está escrito para psicólogos y coordinadores que administran la plataforma sin
necesidad de conocimientos técnicos.

El panel te permite:
\begin{itemize}
  \item Ver y gestionar las solicitudes de citas de los pacientes.
  \item Administrar los usuarios (pacientes y terapeutas).
  \item Publicar noticias, artículos y contenido para el sitio web.
  \item Gestionar la publicidad y las campañas.
  \item Revisar la contabilidad básica.
  \item Recibir notificaciones en tiempo real de lo que pasa en la plataforma.
\end{itemize}

% ================================================================
\section{Cómo ingresar al panel}

\begin{enumerate}
  \item Abre tu navegador (Chrome, Firefox o Edge).
  \item Ve a la dirección del panel (tu administrador te la dará, por ejemplo:\\
        \texttt{https://tudominio.com/admin/login}).
  \item Escribe tu correo electrónico y tu contraseña.
  \item Haz clic en \textbf{Iniciar sesión}.
\end{enumerate}

\begin{importante}
  Si olvidaste tu contraseña, haz clic en \textbf{¿Olvidaste tu contraseña?} en
  la pantalla de login. Recibirás un correo con instrucciones para crear una nueva.
  El enlace del correo expira en 15 minutos.
\end{importante}

% ================================================================
\section{El panel principal}

Al ingresar, verás el \textbf{Panel Principal} con:
\begin{itemize}
  \item Un resumen de actividad reciente.
  \item Accesos directos a las secciones más usadas: Solicitudes, Usuarios,
        Publicidad y Contabilidad.
  \item El menú lateral izquierdo para navegar a cualquier sección.
\end{itemize}

La \textbf{campana} en la esquina superior derecha del menú muestra las
notificaciones en tiempo real. Un número rojo indica cuántas notificaciones
no has leído todavía.

% ================================================================
\section{Notificaciones}

Las notificaciones te avisan cuando ocurre algo importante sin que tengas que
revisar el sistema manualmente.

\subsection{Tipos de notificaciones}

\begin{tabular}{@{}ll@{}}
  \toprule
  \textbf{Notificación} & \textbf{Qué significa} \\
  \midrule
  Nueva solicitud de cita & Un paciente pidió una cita \\
  Cita confirmada & Una cita pasó a estado Confirmada \\
  Cita cancelada & Un paciente o terapeuta canceló \\
  Cita completada & La sesión se marcó como realizada \\
  Paciente no se presentó & El paciente no llegó a la sesión \\
  \bottomrule
\end{tabular}

\subsection{Cómo ver todas las notificaciones}

Haz clic en la campana y luego en \textbf{Ver todas las notificaciones}, o
ve a \textbf{Notificaciones} en el menú lateral.

Desde esa pantalla puedes:
\begin{itemize}
  \item Filtrar solo las no leídas.
  \item Marcar una notificación específica como leída.
  \item Marcar todas como leídas de una vez.
\end{itemize}

% ================================================================
\section{Gestión de citas}

Ve a \textbf{Citas de los pacientes} en el menú lateral.

\subsection{Estados de una cita}

\begin{tabular}{@{}lp{9cm}@{}}
  \toprule
  \textbf{Estado} & \textbf{Significado} \\
  \midrule
  REQUESTED   & El paciente solicitó la cita. Falta confirmarla. \\
  CONFIRMED   & La cita fue aceptada y está programada. \\
  COMPLETED   & La sesión se realizó con éxito. \\
  CANCELLED   & La cita fue cancelada (paciente, terapeuta o admin). \\
  NO\_SHOW    & El paciente no se presentó. \\
  \bottomrule
\end{tabular}

\subsection{Cambiar el estado de una cita}

\begin{enumerate}
  \item Busca la cita en la lista (puedes usar el buscador).
  \item Haz clic en \textbf{Editar}.
  \item Selecciona el nuevo estado en el menú desplegable.
  \item Escribe una nota en \textbf{Notas del administrador} si quieres dejar un comentario.
  \item Haz clic en \textbf{Guardar}.
\end{enumerate}

\begin{importante}
  No puedes leer las \textbf{notas privadas} que el paciente escribió para su terapeuta.
  Esas notas son confidenciales y el sistema las protege automáticamente.
  Las \textbf{Notas del administrador} sí son visibles y editables para ti.
\end{importante}

\subsection{Agendar una cita para un paciente}

Si un paciente te llama por teléfono y quiere reservar:
\begin{enumerate}
  \item Ve a \textbf{Agendar una cita} en el menú.
  \item Selecciona el paciente de la lista.
  \item Elige el terapeuta, el servicio y el horario disponible.
  \item Haz clic en \textbf{Confirmar reserva}.
\end{enumerate}

% ================================================================
\section{Usuarios}

Ve a \textbf{Personas registradas} en el menú lateral.

Aquí puedes ver todos los pacientes y terapeutas registrados. Las acciones
disponibles son:

\begin{itemize}
  \item \textbf{Buscar} por nombre o correo usando el campo de búsqueda.
  \item \textbf{Filtrar} por estado: Activo, Inactivo, Suspendido.
  \item \textbf{Cambiar el estado} de un usuario (por ejemplo, suspender una cuenta).
  \item \textbf{Actualizar el avatar} (foto de perfil) de un terapeuta.
\end{itemize}

\begin{consejo}
  Usa el filtro de estado para encontrar rápidamente cuentas pendientes
  de activación o usuarios suspendidos.
\end{consejo}

% ================================================================
\section{Contenido del sitio web}

Esta sección te permite gestionar todo lo que aparece en el sitio público.

\subsection{Páginas del sitio web}

Ve a \textbf{Páginas del sitio web} para editar páginas como ``Sobre nosotros'',
``Política de privacidad'', etc. Cada página tiene elementos (secciones de texto,
imágenes, botones) que puedes editar individualmente.

\subsection{Noticias y publicaciones}

\begin{enumerate}
  \item Ve a \textbf{Noticias y columnas} para ver el listado de publicaciones.
  \item Haz clic en \textbf{Nueva publicación} para crear un artículo.
  \item Escribe el título, el contenido y selecciona el autor y las categorías.
  \item Elige el estado: \textbf{Borrador} (no visible) o \textbf{Publicado} (visible).
  \item Haz clic en \textbf{Guardar}.
\end{enumerate}

\subsection{Suscriptores premium}

En \textbf{Suscriptores premium} puedes ver quién tiene acceso al contenido
exclusivo y gestionar sus suscripciones.

% ================================================================
\section{Publicidad}

El módulo de publicidad te permite gestionar los anuncios que aparecen en el sitio.

\begin{itemize}
  \item \textbf{Empresas}: Las organizaciones que pagan por publicidad.
  \item \textbf{Dónde se muestra}: Las ubicaciones disponibles en el sitio (banner superior, lateral, etc.).
  \item \textbf{Campañas}: Cada campaña pertenece a una empresa y define fechas y presupuesto.
  \item \textbf{Imágenes}: Los creativos (banners, ilustraciones) de cada campaña.
\end{itemize}

% ================================================================
\section{Productos terapéuticos}

\subsection{Enfoques terapéuticos}
Son los tipos de terapia que ofrece la plataforma (ej: Terapia cognitivo-conductual).

\subsection{Servicios}
Son los paquetes o sesiones concretas que puede reservar un paciente, con precio
y duración definidos. Para desactivar un servicio sin eliminarlo, cambia su estado
a \textbf{Inactivo} — dejará de aparecer en el catálogo público.

% ================================================================
\section{Imágenes y documentos}

Ve a \textbf{Imágenes y documentos} para gestionar los archivos subidos a la plataforma
(fotos de terapeutas, documentos, imágenes del sitio web).

Puedes:
\begin{itemize}
  \item Ver todos los archivos subidos.
  \item Filtrar por tipo (imagen, documento) o por módulo.
  \item Ver a qué entidad está asociado cada archivo.
\end{itemize}

% ================================================================
\section{Contabilidad}

El módulo de contabilidad es para el contador de la plataforma, pero los
administradores pueden consultarlo para revisar transacciones.

\begin{itemize}
  \item \textbf{Cuentas}: El catálogo de cuentas contables.
  \item \textbf{Transacciones}: Movimientos de ingresos y gastos.
  \item \textbf{Pago de citas}: Desde la cita puedes marcarla como pagada,
        lo que genera automáticamente un registro contable.
\end{itemize}

% ================================================================
\section{Preguntas frecuentes}

\subsection{¿Por qué no puedo ver las notas que un paciente le escribió a su terapeuta?}
Las notas entre paciente y terapeuta son confidenciales. El sistema las protege
automáticamente y no las muestra en ninguna pantalla de administración.
Solo el terapeuta puede verlas.

\subsection{¿Qué pasa si dos personas intentan reservar el mismo horario al mismo tiempo?}
El sistema maneja esto automáticamente. Solo una de las dos reservas se completará.
La otra recibirá un mensaje indicando que el horario ya no está disponible.

\subsection{¿Puedo recuperar una cita cancelada?}
No es posible revertir una cancelación directamente. Deberás crear una nueva cita.

\subsection{¿Qué significa el número rojo en la campana?}
Indica cuántas notificaciones nuevas tienes sin leer. Al hacer clic verás un
resumen; para ver el historial completo, ve a la página de Notificaciones.

\subsection{¿Cada cuánto se actualizan las notificaciones?}
En tiempo real, sin necesidad de recargar la página. Cuando ocurre algo nuevo
(por ejemplo, un paciente solicita una cita), la campana se actualiza automáticamente.

% ================================================================
\section{Cierre de sesión}

Para salir del panel de forma segura, haz clic en \textbf{Cerrar sesión} en la
parte inferior del menú lateral. Esto invalida tu sesión activa.

\begin{importante}
  No cierres la ventana del navegador sin cerrar sesión si estás usando un
  computador compartido o público.
\end{importante}

% ================================================================
\section{Soporte técnico}

Si encuentras un error o necesitas ayuda, comunícate con el equipo técnico
indicando:
\begin{itemize}
  \item Qué estabas haciendo cuando ocurrió el problema.
  \item El mensaje de error que apareció (si lo hay).
  \item La fecha y hora aproximadas.
\end{itemize}

\vfill
\begin{center}
  \textcolor{grayText}{\small
    Corazón Migrante \textperiodcentered{} Manual del Panel Administrativo \textperiodcentered{} Versión 1.0 \textperiodcentered{} Julio 2026
  }
\end{center}

\end{document}
